\section{The heat-vector bridge}
\label{sec:heat-vectors}
We now supply the independent-height Gram limit in
Input~\ref{input:word-boundary} from the joint boundary theorem. Throughout
this section the complex heat and slab properties used in
Section~\ref{sec:analytic-boundary} are retained. The involution $J$ is the
one induced on the field union by Euclidean time reflection and conjugation.
This identification is essential in the Gram formula below.

For $0<y_1<\cdots<y_n$, real positions $X_i=(t_i,x_i)$, and an admissible
width $p$ satisfying $4p<y_1$ and $4p<y_j-y_i$ for $i<j$, define
\begin{multline}\label{eq:point-heat-vector}
 Q_{p,y}(X;v)=E(-i(t_1+iy_1)-p)U(x_1)B_p(v_1)\\
 {}\cdot R_p(t_2-t_1+i(y_2-y_1),x_2-x_1)B_p(v_2)\cdots
 R_p(t_n-t_{n-1}+i(y_n-y_{n-1}),x_n-x_{n-1})B_p(v_n)\Omega_E.
\end{multline}
The empty tuple gives $\Omega_E$. The subtraction in the first complex
heat time is $p$, whereas that between two insertions is $2p$.
Each complex heat factor is a contraction, so
\begin{equation}\label{eq:point-heat-bound}
 \|Q_{p,y}(X;v)\|\leq\prod_i\|B_p(v_i)\|.
\end{equation}
For fixed $p,y,v$ this vector is continuous in $X$, by the same successive
strong-continuity argument used for \eqref{eq:continued-matrix}. Consequently
$f(X)Q_{p,y}(X;v)$ is Bochner integrable for every joint Schwartz test $f$.
To justify strong measurability without a global separability assumption on
the Hilbert space, its continuous image of the separable position space is
separable. The bound \eqref{eq:point-heat-bound} then gives the integrable
majorant $|f(X)|\prod_i\|B_p(v_i)\|$.

\begin{lemma}[Mixed Gram identity]\label{lem:mixed-heat-gram}
For positive strictly ordered height tuples $y,z$, separately admissible
widths $p,q$, positions $X,Y$ and labels $v,u$, put
\[
 \widehat y=(-y_n,\ldots,-y_1,z_1,\ldots,z_m),\quad
 \widehat X=(X_n,\ldots,X_1,Y_1,\ldots,Y_m),\quad
 \widehat v=(Jv_n,\ldots,Jv_1,u_1,\ldots,u_m).
\]
Then
\begin{equation}\label{eq:point-gram-continuation}
 \langle Q_{p,y}(X;v),Q_{q,z}(Y;u)\rangle
   =F_{n+m}(\widehat t+i\widehat y,\widehat x;\widehat v).
\end{equation}
Here $\widehat t,\widehat x$ are the real time and spatial parts of
$\widehat X$. The real positions in the first tuple are reversed in order;
their real time coordinates are not negated.
\end{lemma}
\begin{proof}
When all real time coordinates are zero, the vector in
\eqref{eq:point-heat-vector} is the Euclidean word vector with insertions
$(x_i,y_i)$ and labels $v_i$. This follows directly by induction from the
right: a centered insertion adds its point at height $p$, the preceding
positive heat interval supplies the missing distance, and the spatial
translation changes the common reference point. More explicitly, with a
preceding reference $(X_0,Y_0)$ and preceding half-width $q_0$, the partial
vector has points
$(x_i-X_0,y_i-Y_0-q_0)$. The recursion uses first heat time
$y_1-Y_0-q_0-p$; after insertion it produces the same formula with the
first point included. Taking $q_0=0$ and $(X_0,Y_0)=0$ gives the stated
word. Pairing the two Euclidean words reflects the first tuple in time and
applies $J$ to its labels. Permutation symmetry allows its reversed order,
and the imaginary-time identity of Lemma~\ref{lem:continued-matrix} proves
\eqref{eq:point-gram-continuation} in this case. This argument does not
require equal widths $p,q$.

For arbitrary fixed real positions, introduce one complex parameter $c$.
In the first vector replace its times by $iy_i+c t_i$, and in the second
by $iz_j+c s_j$. Call the resulting holomorphic vector functions $A(c)$
and $B(c)$. Since all clearance inequalities are strict and there are only
finitely many of them, there is a $\delta>0$ such that these vectors are
well-defined on $|\operatorname{Im}c|<\delta$. The same strip works for
$A(\overline c)$, after decreasing $\delta$ if necessary. The scalar
function
\[
 c\longmapsto\langle A(\overline c),B(c)\rangle
\]
is holomorphic: expand the difference quotient and use antilinearity in
the first argument, which cancels the conjugation of its increment.
Its proposed right side is $F$ at times
$i\widehat y+c\widehat t$, with the fixed combined spatial tuple.
This too is holomorphic on a sufficiently small horizontal strip; its
heights are the concatenation of the reversed negatives of
$y_i-(\operatorname{Im}c)t_i$ and the tuple
$z_j+(\operatorname{Im}c)s_j$, and remain strictly ordered there.

At $c=i\tau$ near zero, both vectors have purely imaginary times, with
these modified positive height tuples. The already proved Euclidean
identity therefore equates the two scalar functions. The identity theorem
on the connected strip extends equality to $c=1$. This is the desired
identity with arbitrary real positions. Empty tuples cause no difficulty:
their vector is the fixed vacuum and the same argument applies.
\end{proof}

For a word $b=((v_1,f_1),\ldots,(v_n,f_n))$, choose
\[
 y(t)=(t,2t,\ldots,nt),\qquad p(t)=\tfrac18\min(t,1),\quad t>0,
\]
and define
\[
 h_t(b)=\int\Big(\prod_i f_i(X_i)\Big)
                  Q_{p(t),y(t)}(X;v)\,dX.
\]
The clearance conditions hold for every $t>0$, and the fixed-height bound
above proves integrability. No uniform pointwise bound as $t\downarrow0$
is asserted or needed.

\begin{proposition}\label{prop:heat-cups}
These vectors satisfy \eqref{eq:independent-heat-gram} for the boundary
functional of Theorem~\ref{thm:joint-boundary}, and $h_t(1)=\Omega_E$.
Moreover, let $T_E(a)=E(-ia^0)U(\mathbf a)$ be the joint real-time
translation representation on $\mathcal H_E$, and for $a\in\mathbb R^4$
let $\tau_a$ act on words letterwise by replacing each test $f$ by
$f(\,\cdot-a)$. Then for every $a\in\mathbb R^4$, every word $b$ and
every $t>0$,
\[
 h_t(\tau_ab)=T_E(a)h_t(b).
\]
\end{proposition}
\begin{proof}
For words $b,c$ of lengths $n,m$, use
Lemma~\ref{lem:mixed-heat-gram} inside the scalar product of their
Bochner integrals. The integrand is absolutely integrable on the product
space: its absolute value is bounded by the product of the two integrable
scalar test magnitudes times the two fixed-height bounds
\eqref{eq:point-heat-bound}. Fubini and the coordinate permutation
$(X,Y)\mapsto(X_n,\ldots,X_1,Y_1,\ldots,Y_m)$ therefore give
\[
 \langle h_s(b),h_t(c)\rangle
 =\psi_{(-ns,\ldots,-s,t,\ldots,mt)}(f_{b^*c}),
\]
with labels $(Jv_n,\ldots,Jv_1,u_1,\ldots,u_m)$ and joint product test
$f_{b^*c}$, whose first factors are the reversed conjugates of the tests
in $b$. The displayed heights are strictly ordered for every $s,t>0$
and approach zero as the two parameters independently approach zero.
Theorem~\ref{thm:joint-boundary} gives the limit $w(b^*c)=K(b,c)$.
The zero-dimensional integral for the empty word is the vacuum itself.

For translations write $T_E(a)=E(-ia^0)U(\mathbf a)$, where the
imaginary-axis values of $E$ denote its strongly continuous unitary boundary
values. These boundary values and their law with the interior extension are
part of the complex heat extension used here. Translating all real
positions by $a$ leaves all differences in
\eqref{eq:point-heat-vector} unchanged; its first factor changes by the
commuting factor $T_E(a)$. Hence
\[
 Q_{p,y}(X+a;v)=T_E(a)Q_{p,y}(X;v).
\]
For the empty word the same assertion follows from translation invariance
of the vacuum. Substitute $X\mapsto X+a$ in the integral with translated
tests $f_i(\,\cdot-a)$ and move the bounded operator $T_E(a)$ outside the
Bochner integral. This proves the translation identity for every $t>0$.
\end{proof}

Thus the heat-vector limit required for positivity and for the isometric
embedding of Proposition~\ref{prop:word-motion} follows from the slab
bounds and complex heat extension. Lorentz covariance, spacelike locality,
and retention of nonzero labels at the boundary are distinct remaining
steps; they have not been inferred from the Gram identity.
